\chapter{Conclusões}\label{ch:conclusiones}

Esta tese perguntou como evolui espacialmente a distribuição de vento e precipitação nos ciclones extratropicais do Atlântico Sul, e quais configurações caracterizam essa organização nas fases incipiente, intensificação, maturidade e decaimento. A resposta depende da intensidade do sistema. Na População Global, ambos os campos compartilham o flanco oriental do ciclone durante as fases iniciais e depois se dispersam sem se reorganizar. No subconjunto de maior vorticidade (p90), por sua vez, a partir da maturidade os máximos de vento se concentram ao noroeste e os de precipitação a leste e sudeste, em uma estrutura em arco afastada do centro, com separações superiores a $8^{\circ}$ nas três regiões de gênese.

A hipótese formulada se confirma no essencial. As assimetrias de ambos os campos variam entre fases e entre estratos de intensidade. Entre regiões, a direção da separação se mantém e o que varia é sua magnitude. A segunda parte da hipótese, de que esses padrões não são capturáveis mediante métricas escalares de intensidade central, se cumpre com clareza para a precipitação, cuja relação com a vorticidade deixa de ser significativa na maturidade dos sistemas p90 de SBR e LPB. Para o vento, cumpre-se apenas em parte, já que seu modo dominante permanece correlacionado com a vorticidade em p90, embora esse modo já não descreva por si só a variabilidade do campo. As seções seguintes apresentam os resultados que sustentam essa resposta, na ordem dos objetivos específicos.

\section{Base de dados semilagrangeana (Objetivo Específico 1)}

Construiu-se uma base de dados semilagrangeana que associa os campos horários de vento a 10~m e precipitação total do ERA5 com as trajetórias de \citet{gramcianinov2020analysis} e com a classificação objetiva de fases do \textit{CycloPhaser} \citep{coutodesouza2024new, deSouza2025CycloPhaser}, conservando unicamente os sistemas com ciclo de vida completo. A amostra foi estratificada por região de gênese (SBR, LPB e ARG) e por intensidade dinâmica ($\zeta_{850}^{\min}$), na População Global e no subconjunto p90. A classificação de fases é coerente com a dinâmica dos sistemas, já que 83\% dos ciclones p90 e 72\% da População Global alcançam seu mínimo de vorticidade durante a maturidade. Essa coerência permite comparar os padrões espaciais fase por fase sem que a atribuição de fases introduza uma defasagem sistemática.

\section{Magnitude e localização dos máximos (Objetivo Específico 2)}

Em magnitude, o vento separa com clareza ambas as populações à medida que os sistemas maturam. Na maturidade, a moda dos máximos de vento de p90 situa-se entre 23 e 25~m\,s$^{-1}$, entre 25\% e 50\% acima da População Global (16--19~m\,s$^{-1}$), com caudas de até 30~m\,s$^{-1}$. A precipitação se comporta de outra forma. Nas fases iniciais, p90 se desloca apenas levemente para valores maiores (${\sim}25$~mm\,h$^{-1}$ em SBR incipiente e ${\sim}30$~mm\,h$^{-1}$ em LPB intensificação), e as taxas superiores a 60~mm\,h$^{-1}$ não são exclusivas de p90. Na maturidade a relação se inverte, e a moda de p90 em SBR (${\sim}7$~mm\,h$^{-1}$) fica abaixo da População Global (${\sim}17$~mm\,h$^{-1}$), até se concentrar em 1--4~mm\,h$^{-1}$ no decaimento de SBR e LPB. ARG apresenta as menores taxas de precipitação em todas as fases. A partir da maturidade, portanto, pertencer a p90 implica ventos mais intensos, mas não maior precipitação.

Em localização, a População Global mostra distribuições amplas e de gradientes suaves. Os máximos de vento passam do setor leste nas primeiras fases para o quadrante noroeste na maturidade, e os de precipitação mantêm uma assimetria persistente para leste e sudeste, na posição que a literatura atribui ao WCB. Em p90 essa organização se torna mais compacta. Na maturidade, os máximos de vento se concentram no quadrante noroeste, a ${\sim}2$--$4^{\circ}$ do centro, com densidades de 19--27$\times10^{-4}$ frente a 11--15$\times10^{-4}$ na População Global, na zona onde a literatura situa a esteira transportadora fria (CCB) e a fratura frontal. Os máximos de precipitação, por sua vez, se afastam do centro e se confinam em faixas alongadas a leste e sudeste, que persistem no decaimento enquanto na População Global o sinal se dissipa. A concentração do vento é maior em SBR e LPB do que em ARG, embora SBR apresente a distribuição de magnitudes mais larga, de modo que a região que mais concentra a localização de seus máximos não é a que mais concentra sua intensidade.

\section{Modos dominantes de variabilidade (Objetivo Específico 3)}

A análise EOF univariada mostra que a intensidade muda a forma como se distribui a variabilidade de cada campo, e que o faz em sentidos opostos. Na População Global, o EOF1 de vento explica entre 30.9\% e 57.4\% da variância, é significativo em todo o domínio e sua PC1 se correlaciona com a vorticidade até $|r|=0.86$--$0.89$ na maturidade. Em p90, o EOF1 de vento não supera 37.6\% e a variância se distribui entre vários modos. Sua PC1 permanece correlacionada de forma significativa com a vorticidade ($|r|$ entre $0.26$ e $0.74$), mas PC2 e PC3 superam PC1 em quatro combinações de região e fase (até $|r|=0.58$ em SBR), o que é compatível com um maior peso da variabilidade de mesoescala nos ciclones intensos. A precipitação segue o caminho inverso. Na População Global sua variabilidade se dispersa entre múltiplos modos (EOF1 entre 11.1\% e 16.4\%), enquanto em p90 o EOF1 se concentra na maturidade e no decaimento, até 36.5\% em LPB e 26.2\% em SBR durante o decaimento. Nessas mesmas fases, a correlação de sua PC1 com a vorticidade deixa de ser significativa em SBR ($|r|=0.17$) e LPB ($|r|=0.03$) durante a maturidade e só se mantém em ARG ($|r|=0.24$).

A análise multivariada (MEOF) avalia se essas transformações se refletem na covariabilidade conjunta de ambos os campos. Na População Global, o MEOF1 captura entre 14.8\% e 26.6\% da covariância conjunta, com gradientes de precipitação e vento orientados em eixos distintos (oeste-leste e norte-sul, respectivamente). Em p90, essa fração cai para 11.65\%--23.47\% e, na maturidade, o MEOF1 mostra um padrão de anticovariância com as cargas de precipitação a sudeste e as de vento a oeste. Esse resultado tem uma ressalva importante. O vento domina o MEOF1 na maioria das combinações de região, fase e população, e a precipitação só contribui com força comparável em p90 a partir da maturidade. O padrão de anticovariância é, portanto, interpretável sobretudo nessa combinação. Ali, a correlação significativa entre as PC1 univariadas de ambos os campos ($|r|$ entre $0.31$ e $0.45$ nas três regiões) é consistente com a separação geométrica, já que os eventos que mais projetam sobre a estrutura em arco de precipitação tendem a ser os que menos projetam sobre o núcleo de vento.

\section{Protótipos estruturais (Objetivo Específico 4)}

Os protótipos estruturais, que sobrepõem em um mesmo painel as PDFe e o EOF1 univariado de ambas as variáveis, convertem os resultados anteriores em distâncias concretas entre os máximos de cada campo. Na População Global, as separações modais durante as fases iniciais são pequenas em SBR (abaixo de $1^{\circ}$) e moderadas em LPB e ARG (até $4{,}9^{\circ}$ e $5{,}5^{\circ}$, respectivamente), e na maturidade e no decaimento os contornos permanecem extensos, sem que nenhum campo adote uma geometria bem definida, embora em ARG as modas já se separem ${\sim}10^{\circ}$. Em p90, na maturidade de SBR, a moda de precipitação ($\Delta x = 6{,}8^{\circ}$, $\Delta y = -1{,}9^{\circ}$) e a de vento ($\Delta x = -1{,}8^{\circ}$, $\Delta y = 2{,}4^{\circ}$) ficam separadas ${\sim}9{,}5^{\circ}$, distância que se mantém no decaimento. LPB apresenta separações um pouco menores (${\sim}8$--$9^{\circ}$) e ARG as maiores (${\sim}10$--$11^{\circ}$).

A direção da separação, precipitação a leste-sudeste e vento a noroeste, se repete nas três regiões, enquanto sua magnitude varia sem que esta tese identifique um único fator regional que a explique. Em SBR e LPB, a umidade canalizada pelo Jato de Camadas Baixas da América do Sul (SALLJ) poderia sustentar a estrutura em arco com maior intensidade. ARG, onde a correlação entre o PC1 de precipitação e a vorticidade se mantém significativa durante todo o ciclo de vida em p90, apresenta no entanto a maior separação das três regiões. Essa configuração é compatível com a transição ao modelo de Shapiro-Keyser, na qual uma possível intrusão seca separaria ambos os campos, embora esse mecanismo não seja medido diretamente nesta tese e sua interpretação se apoie em estudos do Hemisfério Norte.

Em síntese, em resposta ao objetivo geral, a estrutura espacial do vento e da precipitação nos ciclones extratropicais do Atlântico Sul evolui de forma distinta segundo a intensidade do sistema. Nos ciclones de maior vorticidade, os máximos de ambos os campos terminam em setores opostos durante a maturidade e o decaimento, uma configuração que a média da População Global não mostra e que a vorticidade do centro, por si só, não descreve, o que é relevante para o risco costeiro.
